\documentclass[aspectratio=169,11pt]{beamer}

% ------------------------------------------------------------
% Tema y paquetes
% ------------------------------------------------------------
\usetheme{Madrid}
\usecolortheme{seahorse}
\setbeamertemplate{navigation symbols}{}
\setbeamertemplate{footline}[frame number]
\usepackage[spanish,es-nodecimaldot]{babel}
\usepackage[utf8]{inputenc}
\usepackage[T1]{fontenc}
\usepackage{lmodern}
\usepackage{amsmath, amssymb, bm}
\usepackage{booktabs}
\usepackage{array}
\usepackage{multicol}
\usepackage{tikz}
\usepackage{listings}
\usepackage{xcolor}

% ------------------------------------------------------------
% Colores
% ------------------------------------------------------------
\definecolor{macropblue}{RGB}{0,64,128}
\definecolor{macropgray}{RGB}{80,80,80}
\definecolor{macropred}{RGB}{150,30,30}
\definecolor{macropgreen}{RGB}{20,120,80}
\definecolor{codebg}{RGB}{245,247,250}

\setbeamercolor{title}{fg=white,bg=macropblue}
\setbeamercolor{frametitle}{fg=white,bg=macropblue}
\setbeamercolor{block title}{fg=white,bg=macropblue}
\setbeamercolor{block body}{bg=blue!3}
\setbeamercolor{alerted text}{fg=macropred}

% ------------------------------------------------------------
% Código R
% ------------------------------------------------------------
\lstdefinelanguage{R}{
  keywords={library, data.frame, lm, summary, predict, mean, sd, quantile, cbind, matrix,
  if, else, for, in, function, return, list, rep, seq, set.seed, rnorm, qnorm, pnorm},
  otherkeywords={<-, ~, $, \%, \%*\%},
  sensitive=true,
  morecomment=[l]{\#},
  morestring=[b]",
  morestring=[b]'
}
\lstdefinestyle{Rstyle}{
  language=R,
  backgroundcolor=\color{codebg},
  basicstyle=\ttfamily\scriptsize,
  keywordstyle=\color{macropblue}\bfseries,
  commentstyle=\color{macropgreen},
  stringstyle=\color{macropred},
  frame=single,
  breaklines=true,
  showstringspaces=false,
  columns=fullflexible,
  literate={á}{{\'a}}1 {é}{{\'e}}1 {í}{{\'i}}1 {ó}{{\'o}}1 {ú}{{\'u}}1
           {Á}{{\'A}}1 {É}{{\'E}}1 {Í}{{\'I}}1 {Ó}{{\'O}}1 {Ú}{{\'U}}1
           {ñ}{{\~n}}1 {Ñ}{{\~N}}1 {ü}{{\"u}}1 {¿}{{?`}}1 {¡}{{!`}}1
}

% ------------------------------------------------------------
% Metadatos
% ------------------------------------------------------------
\title[Stress Testing del Portafolio]{Módulo 4: Stress Testing del Portafolio Previsional}
\subtitle{De escenarios macro-financieros a pérdidas, rentabilidad y suficiencia previsional}
\author[MACROPOL]{MACROPOL: Economía Aplicada y Políticas Públicas}
\institute{Curso: Stress Testing en Sistemas de Pensiones de Capitalización Individual}
\date{2026}

% ------------------------------------------------------------
\begin{document}

% ============================================================
\begin{frame}
  \titlepage
\end{frame}

% ============================================================
\begin{frame}{Ruta del módulo}
\begin{block}{Objetivo}
Traducir los escenarios macro-financieros construidos en el Módulo 3 en impactos cuantitativos sobre el portafolio previsional, la rentabilidad acumulada y la suficiencia futura de las pensiones.
\end{block}

\begin{enumerate}
  \item Conexión con Módulo 3: escenario, narrativa y trayectorias.
  \item Modelo multifactorial del portafolio previsional.
  \item Volatilidad condicional: ARCH/GARCH.
  \item Correlaciones en crisis y pérdida de diversificación.
  \item Simulación de rentabilidad y pérdidas extremas.
  \item Lectura regulatoria: límites, concentración y resiliencia.
\end{enumerate}
\end{frame}

% ============================================================
\begin{frame}{Conexión directa con el Módulo 3}
\begin{columns}[T]
\column{0.48\textwidth}
\begin{block}{Módulo 3: Construcción de escenarios}
\begin{itemize}
  \item ARIMA para variables individuales.
  \item VAR para interdependencias macro-financieras.
  \item Monte Carlo para trayectorias alternativas.
  \item Choques históricos vs. hipotéticos.
\end{itemize}
\end{block}

\column{0.48\textwidth}
\begin{block}{Módulo 4: Impacto en portafolio}
\begin{itemize}
  \item Betas del portafolio frente a factores.
  \item Volatilidad condicional bajo estrés.
  \item Correlaciones que aumentan en crisis.
  \item Distribución de rentabilidades y pérdidas.
\end{itemize}
\end{block}
\end{columns}

\vspace{0.4cm}
\centering
\Large
\alert{El escenario macro-financiero deja de ser una narrativa: se convierte en pérdida previsional cuantificable.}
\end{frame}

% ============================================================
\begin{frame}{Del Módulo 3 al Módulo 4: insumos y productos}
\begin{center}
\begin{tikzpicture}[node distance=1.45cm, every node/.style={font=\small}]
\node[draw, rounded corners, fill=blue!8, text width=3.0cm, align=center] (m3) {Módulo 3\\ Escenario macro-financiero};
\node[draw, rounded corners, fill=blue!8, right of=m3, xshift=2.6cm, text width=3.0cm, align=center] (factores) {Factores de riesgo\\ tasas, FX, equity, spreads};
\node[draw, rounded corners, fill=blue!8, right of=factores, xshift=2.8cm, text width=3.0cm, align=center] (port) {Modelo de portafolio\\ sensibilidades y pesos};
\node[draw, rounded corners, fill=red!8, right of=port, xshift=2.6cm, text width=2.0cm, align=center] (impacto) {Pérdidas\\ rentabilidad y saldo};

\draw[->, thick] (m3) -- (factores);
\draw[->, thick] (factores) -- (port);
\draw[->, thick] (port) -- (impacto);
\end{tikzpicture}
\end{center}

\begin{block}{Producto del Módulo 4}
Una matriz de impacto que conecte cada choque macro-financiero con:
\begin{itemize}
  \item pérdida de valor del portafolio,
  \item rentabilidad acumulada bajo escenario adverso,
  \item sensibilidad por clase de activo,
  \item implicación sobre tasas de reemplazo futuras.
\end{itemize}
\end{block}
\end{frame}

% ============================================================
\begin{frame}{Baseline vs. adverso: separación conceptual}
\begin{columns}[T]
\column{0.48\textwidth}
\begin{block}{Escenario baseline}
\begin{itemize}
  \item Crecimiento converge a potencial.
  \item Inflación retorna a meta.
  \item Tasas se normalizan.
  \item Spreads financieros estables.
  \item Correlaciones históricas promedio.
\end{itemize}
\end{block}

\column{0.48\textwidth}
\begin{alertblock}{Escenario adverso}
\begin{itemize}
  \item Recesión o desaceleración fuerte.
  \item Inflación persistente.
  \item Aumento de tasas y spreads.
  \item Depreciación cambiaria.
  \item Correlaciones aumentan y la diversificación falla.
\end{itemize}
\end{alertblock}
\end{columns}

\vspace{0.3cm}
\begin{block}{Regla de enseñanza}
Baseline describe la trayectoria esperada. El adverso evalúa resiliencia, no predicción.
\end{block}
\end{frame}

% ============================================================
\begin{frame}{Narrativa económica vs. modelo cuantitativo}
\begin{table}
\centering
\small
\begin{tabular}{p{0.45\textwidth} p{0.45\textwidth}}
\toprule
\textbf{Narrativa económica} & \textbf{Modelo cuantitativo} \\
\midrule
``Shock externo eleva tasas internacionales y reduce apetito por riesgo'' &
Trayectoria de tasa libre de riesgo, spread soberano, equity global y tipo de cambio. \\
\addlinespace
``Inflación persistente erosiona retornos reales'' &
Deflactación de rentabilidad nominal y ajuste de curva de tasas. \\
\addlinespace
``Los activos se vuelven más correlacionados'' &
Matriz de correlación estresada o modelo DCC-GARCH. \\
\bottomrule
\end{tabular}
\end{table}

\begin{block}{Mensaje central}
La narrativa justifica la plausibilidad del shock; el modelo cuantitativo mide su magnitud en el portafolio.
\end{block}
\end{frame}

% ============================================================
\begin{frame}{Supuestos vs. resultados}
\begin{columns}[T]
\column{0.48\textwidth}
\begin{block}{Supuestos}
\begin{itemize}
  \item Pesos iniciales del portafolio.
  \item Duración de bonos.
  \item Betas frente a factores.
  \item Matriz de correlaciones.
  \item Horizonte del stress test.
\end{itemize}
\end{block}

\column{0.48\textwidth}
\begin{block}{Resultados}
\begin{itemize}
  \item Pérdida total del fondo.
  \item Contribución por factor.
  \item VaR y Expected Shortfall.
  \item Máximo drawdown.
  \item Reducción del saldo acumulado.
\end{itemize}
\end{block}
\end{columns}

\vspace{0.3cm}
\centering
\alert{Una buena presentación regulatoria nunca mezcla supuestos con resultados.}
\end{frame}

% ============================================================
\section{Contexto}

\begin{frame}{Contexto: riesgo financiero como riesgo previsional}
\begin{block}{Intuición económica}
En un sistema de capitalización individual, el afiliado acumula un saldo que depende de aportes, comisiones y rentabilidad. Por tanto, un shock de mercado no solo reduce riqueza financiera: reduce potencialmente la pensión futura.
\end{block}

\[
S_T=\sum_{t=1}^{T} c_t \prod_{j=t}^{T}(1+r_j)
\]

\begin{itemize}
  \item $S_T$: saldo acumulado al retiro.
  \item $c_t$: contribución neta.
  \item $r_j$: rentabilidad del fondo.
\end{itemize}

\begin{alertblock}{Implicación}
La variable crítica del Módulo 4 es $r_t$: la rentabilidad bajo estrés.
\end{alertblock}
\end{frame}

% ============================================================
\begin{frame}{Riesgos principales del portafolio previsional}
\begin{table}
\centering
\small
\begin{tabular}{p{0.27\textwidth} p{0.38\textwidth} p{0.25\textwidth}}
\toprule
\textbf{Riesgo} & \textbf{Canal de pérdida} & \textbf{Indicador} \\
\midrule
Tasa de interés & Caída de precio de bonos ante aumento de tasas & Duración / DV01 \\
Accionario & Menor precio de acciones locales y externas & Beta equity \\
Cambiario & Pérdida o ganancia por exposición en moneda extranjera & Exposición FX \\
Crédito & Aumento de spreads o deterioro de rating & Spread duration \\
Inflación & Reducción de rentabilidad real & Retorno real \\
Liquidez & Venta forzada con descuentos & Haircut / bid-ask \\
Correlación & Falla de diversificación en crisis & Matriz $\rho$ \\
\bottomrule
\end{tabular}
\end{table}
\end{frame}

% ============================================================
\begin{frame}{Mapa de transmisión macro-financiera}
\begin{center}
\begin{tikzpicture}[node distance=1.35cm, every node/.style={font=\small}]
\node[draw, rounded corners, fill=blue!8, text width=4.1cm, align=center] (shock) {Shock macro\\ inflación, crecimiento, tasas};
\node[draw, rounded corners, fill=blue!8, below of=shock, text width=5.1cm, align=center] (mercado) {Precios financieros\\ bonos, acciones, FX, spreads};
\node[draw, rounded corners, fill=blue!8, below of=mercado, text width=4.1cm, align=center] (rent) {Rentabilidad del fondo\\ nominal y real};
\node[draw, rounded corners, fill=red!8, below of=rent, text width=3.1cm, align=center] (saldo) {Saldo individual\\ acumulado};
\node[draw, rounded corners, fill=red!8, below of=saldo, text width=3.1cm, align=center] (pension) {Pensión esperada\\ tasa de reemplazo};

\draw[->, thick] (shock) -- (mercado);
\draw[->, thick] (mercado) -- (rent);
\draw[->, thick] (rent) -- (saldo);
\draw[->, thick] (saldo) -- (pension);
\end{tikzpicture}
\end{center}
\end{frame}

% ============================================================
\begin{frame}{Pregunta regulatoria del módulo}
\begin{block}{Pregunta central}
¿Bajo qué combinaciones de choques macro-financieros el portafolio previsional pierde valor suficiente para comprometer la suficiencia futura de las pensiones?
\end{block}

\begin{columns}[T]
\column{0.32\textwidth}
\begin{block}{Micro}
Impacto por fondo, clase de activo y cohorte de afiliados.
\end{block}

\column{0.32\textwidth}
\begin{block}{Macro}
Impacto agregado sobre ahorro previsional y confianza del sistema.
\end{block}

\column{0.32\textwidth}
\begin{block}{Política}
Necesidad de límites, coberturas, liquidez y reglas de ciclo de vida.
\end{block}
\end{columns}
\end{frame}

% ============================================================
\section{Modelo / Herramienta}

\begin{frame}{Modelo multifactorial del portafolio}
\begin{block}{Modelo base}
\[
R_{p,t}=\alpha+\bm{\beta}'\bm{F}_t+\varepsilon_t
\]
\end{block}

\begin{itemize}
  \item $R_{p,t}$: retorno del portafolio previsional.
  \item $\bm{F}_t$: vector de factores macro-financieros.
  \item $\bm{\beta}$: sensibilidad del fondo a cada factor.
  \item $\varepsilon_t$: componente idiosincrático no explicado.
\end{itemize}

\begin{alertblock}{Lectura previsional}
Las betas muestran qué factores del escenario adverso reducen más la rentabilidad del fondo.
\end{alertblock}
\end{frame}

% ============================================================
\begin{frame}{Factores del Módulo 3 que entran al Módulo 4}
\begin{table}
\centering
\small
\begin{tabular}{p{0.30\textwidth} p{0.33\textwidth} p{0.25\textwidth}}
\toprule
\textbf{Variable del escenario} & \textbf{Factor financiero} & \textbf{Activo afectado} \\
\midrule
Tasa de política & Curva soberana local & Bonos locales \\
Inflación esperada & Tasa real / retorno real & Bonos nominales y saldo real \\
Tipo de cambio & Retorno FX & Activos externos \\
PIB / actividad & Prima de riesgo accionario & Acciones locales \\
Tasa Fed / global & Curva internacional & Bonos externos \\
Riesgo país & Spread soberano & Bonos soberanos y corporativos \\
\bottomrule
\end{tabular}
\end{table}
\end{frame}

% ============================================================
\begin{frame}{Retorno del portafolio por clases de activo}
\[
R_{p,t}=\sum_{i=1}^{N} w_i R_{i,t}
\]

\begin{block}{Conexión con escenarios}
Cada $R_{i,t}$ puede modelarse como función de factores:
\[
R_{i,t}=a_i+\bm{b}_i'\bm{F}_t+u_{i,t}
\]
\end{block}

\begin{itemize}
  \item Los pesos $w_i$ reflejan la política de inversión.
  \item Las sensibilidades $\bm{b}_i$ reflejan riesgo de mercado.
  \item El escenario adverso genera una trayectoria para $\bm{F}_t$.
\end{itemize}
\end{frame}

% ============================================================
\begin{frame}{Renta fija: duración y choque de tasas}
\begin{block}{Aproximación de primer orden}
\[
\frac{\Delta P}{P}\approx -D_{\text{mod}}\Delta y
\]
\end{block}

\begin{itemize}
  \item $D_{\text{mod}}$: duración modificada.
  \item $\Delta y$: cambio en rendimiento.
  \item Un aumento de tasas genera pérdida de precio.
\end{itemize}

\begin{alertblock}{Implicación para pensiones}
Fondos con alta proporción de bonos largos son vulnerables a escenarios de inflación persistente y endurecimiento monetario.
\end{alertblock}
\end{frame}

% ============================================================
\begin{frame}{Renta fija: duración, convexidad y estrés}
\begin{block}{Aproximación con convexidad}
\[
\frac{\Delta P}{P}\approx -D_{\text{mod}}\Delta y+\frac{1}{2}C(\Delta y)^2
\]
\end{block}

\begin{itemize}
  \item La convexidad corrige la pérdida para choques grandes.
  \item En stress testing, choques de tasas suelen ser suficientemente grandes para requerir convexidad.
  \item Para fondos previsionales, importa el efecto acumulado sobre el saldo, no solo la pérdida instantánea.
\end{itemize}
\end{frame}

% ============================================================
\begin{frame}{Acciones y activos externos}
\begin{block}{Modelo de retorno accionario}
\[
R^{eq}_{t}=\alpha+\beta_g R^{global}_{t}+\beta_l R^{local}_{t}+\beta_{fx}\Delta e_t+\varepsilon_t
\]
\end{block}

\begin{itemize}
  \item $\beta_g$: sensibilidad al ciclo financiero global.
  \item $\beta_l$: sensibilidad al mercado local.
  \item $\beta_{fx}$: exposición cambiaria neta.
\end{itemize}

\begin{alertblock}{Lectura de estrés}
Una depreciación puede compensar pérdidas externas en moneda local, pero aumenta riesgo si hay descalce o cobertura incompleta.
\end{alertblock}
\end{frame}

% ============================================================
\begin{frame}{Volatilidad condicional: por qué GARCH}
\begin{block}{Hecho estilizado}
Los retornos financieros presentan \textit{volatility clustering}: períodos tranquilos y períodos turbulentos.
\end{block}

\begin{columns}[T]
\column{0.48\textwidth}
\begin{block}{Volatilidad constante}
\[
R_t=\mu+\sigma z_t
\]
Supone que el riesgo es igual todos los días.
\end{block}

\column{0.48\textwidth}
\begin{alertblock}{Volatilidad condicional}
\[
R_t=\mu+\sigma_t z_t
\]
Permite que el riesgo aumente durante crisis.
\end{alertblock}
\end{columns}
\end{frame}

% ============================================================
\begin{frame}{ARCH y GARCH}
\begin{block}{ARCH(1)}
\[
\sigma_t^2=\omega+\alpha \varepsilon_{t-1}^2
\]
\end{block}

\begin{block}{GARCH(1,1)}
\[
\sigma_t^2=\omega+\alpha \varepsilon_{t-1}^2+\beta\sigma_{t-1}^2
\]
\end{block}

\begin{table}
\centering
\small
\begin{tabular}{lll}
\toprule
\textbf{Parámetro} & \textbf{Lectura técnica} & \textbf{Lectura económica} \\
\midrule
$\alpha$ & respuesta a shocks recientes & intensidad del pánico \\
$\beta$ & persistencia de volatilidad & duración de la crisis \\
$\alpha+\beta$ & persistencia total & velocidad de normalización \\
\bottomrule
\end{tabular}
\end{table}
\end{frame}

% ============================================================
\begin{frame}{Correlaciones bajo estrés}
\begin{block}{Problema}
La diversificación funciona en tiempos normales, pero puede debilitarse cuando los activos caen simultáneamente.
\end{block}

\[
\sigma^2_{p,t}=\bm{w}'\bm{\Sigma}_t\bm{w}
\]

\[
\bm{\Sigma}_t=
\begin{bmatrix}
\sigma_{1,t}^2 & \rho_{12,t}\sigma_{1,t}\sigma_{2,t} \\
\rho_{21,t}\sigma_{2,t}\sigma_{1,t} & \sigma_{2,t}^2
\end{bmatrix}
\]

\begin{alertblock}{Implicación}
Un escenario adverso debe estresar retornos, volatilidades y correlaciones.
\end{alertblock}
\end{frame}

% ============================================================
\begin{frame}{DCC-GARCH: idea general}
\begin{block}{Modelo}
\[
H_t=D_t R_t D_t
\]
\end{block}

\begin{itemize}
  \item $H_t$: matriz de covarianzas condicionales.
  \item $D_t$: matriz diagonal de volatilidades condicionales.
  \item $R_t$: matriz de correlaciones condicionales.
\end{itemize}

\begin{block}{Uso en stress testing}
Permite evaluar escenarios donde no solo aumenta la volatilidad, sino también la co-movilidad entre clases de activo.
\end{block}
\end{frame}

% ============================================================
\begin{frame}{VaR y Expected Shortfall}
\begin{block}{Value-at-Risk}
\[
VaR_{\alpha}=\inf\{l: P(L\leq l)\geq \alpha\}
\]
\end{block}

\begin{block}{Expected Shortfall}
\[
ES_{\alpha}=E[L \mid L>VaR_{\alpha}]
\]
\end{block}

\begin{alertblock}{Lectura regulatoria}
El VaR responde: ``¿cuánto se pierde en un percentil?'' El ES responde: ``¿cuánto se pierde si se materializa la cola?''
\end{alertblock}
\end{frame}

% ============================================================
\begin{frame}{Simulación de rentabilidad bajo escenarios}
\begin{block}{Trayectoria simulada}
\[
R_{p,t}^{(s)}=\alpha+\bm{\beta}'\bm{F}_{t}^{(s)}+\sigma_t z_t^{(s)}
\]
\end{block}

\begin{itemize}
  \item $s$: simulación o trayectoria.
  \item $\bm{F}_{t}^{(s)}$: trayectoria macro-financiera del Módulo 3.
  \item $\sigma_t$: volatilidad condicional estimada en el Módulo 4.
  \item $z_t^{(s)}$: shock aleatorio, normal o t-Student.
\end{itemize}
\end{frame}

% ============================================================
\section{Implementación}

\begin{frame}[fragile]{Estructura de datos en R}
\begin{lstlisting}[style=Rstyle]
# Base mínima para el módulo 4
# ret_port: rentabilidad del fondo previsional
# tasa, fx, equity, spread, inflacion: factores del módulo 3

datos <- data.frame(
  fecha      = fechas,
  ret_port   = ret_fondo,
  tasa       = delta_tasa,
  fx         = delta_tipo_cambio,
  equity     = ret_equity_global,
  spread     = delta_spread_soberano,
  inflacion  = inflacion_mensual
)
\end{lstlisting}

\begin{block}{Puente con Módulo 3}
Las variables \texttt{tasa}, \texttt{fx}, \texttt{equity}, \texttt{spread} e \texttt{inflacion} provienen de las trayectorias baseline y adversas generadas previamente.
\end{block}
\end{frame}

% ============================================================
\begin{frame}[fragile]{Modelo multifactorial en R}
\begin{lstlisting}[style=Rstyle]
modelo_mf <- lm(
  ret_port ~ tasa + fx + equity + spread + inflacion,
  data = datos
)

summary(modelo_mf)
betas <- coef(modelo_mf)
\end{lstlisting}

\begin{block}{Interpretación}
\begin{itemize}
  \item Coeficiente positivo: el factor aumenta la rentabilidad.
  \item Coeficiente negativo: el factor reduce la rentabilidad.
  \item Magnitud: sensibilidad del fondo ante el shock.
\end{itemize}
\end{block}
\end{frame}

% ============================================================
\begin{frame}[fragile]{Aplicación de escenarios baseline y adverso}
\begin{lstlisting}[style=Rstyle]
# escenarios_m3 contiene trayectorias generadas en el módulo 3
# columnas: tasa, fx, equity, spread, inflacion

escenarios_m3$rp_hat <- predict(
  modelo_mf,
  newdata = escenarios_m3
)

aggregate(rp_hat ~ escenario, data = escenarios_m3, mean)
\end{lstlisting}

\begin{alertblock}{Resultado esperado}
El mismo modelo se aplica a dos narrativas: baseline y adversa. La diferencia entre ambos mide el costo financiero del estrés.
\end{alertblock}
\end{frame}

% ============================================================
\begin{frame}[fragile]{GARCH en R}
\begin{lstlisting}[style=Rstyle]
library(rugarch)

spec <- ugarchspec(
  variance.model = list(model = "sGARCH",
                        garchOrder = c(1, 1)),
  mean.model     = list(armaOrder = c(0, 0)),
  distribution.model = "std"
)

fit_garch <- ugarchfit(spec = spec, data = datos$ret_port)
show(fit_garch)
\end{lstlisting}

\begin{block}{Uso}
La volatilidad condicional estimada permite proyectar pérdidas más realistas durante episodios de estrés.
\end{block}
\end{frame}

% ============================================================
\begin{frame}[fragile]{Volatilidad condicional y VaR en R}
\begin{lstlisting}[style=Rstyle]
sigma_t <- sigma(fit_garch)
mu_t    <- fitted(fit_garch)

alpha <- 0.99
q_t <- qt(alpha, df = 6)   # aproximación t-Student

VaR_99 <- -(mu_t + sigma_t * q_t)
summary(VaR_99)
\end{lstlisting}

\begin{alertblock}{Cuidado}
El VaR debe interpretarse como pérdida potencial bajo supuestos de distribución. En crisis, conviene complementarlo con Expected Shortfall y escenarios determinísticos.
\end{alertblock}
\end{frame}

% ============================================================
\begin{frame}[fragile]{Simulación Monte Carlo del portafolio}
\begin{lstlisting}[style=Rstyle]
set.seed(123)

B <- 10000
mu <- mean(datos$ret_port)
sig <- sd(datos$ret_port)

sim_ret <- rnorm(B, mean = mu, sd = sig)
loss <- -sim_ret

VaR_95 <- quantile(loss, 0.95)
ES_95  <- mean(loss[loss > VaR_95])

c(VaR_95 = VaR_95, ES_95 = ES_95)
\end{lstlisting}

\begin{block}{Conexión con Módulo 3}
En lugar de simular solo retornos históricos, puede alimentarse la simulación con trayectorias de factores generadas por VAR o ARIMA.
\end{block}
\end{frame}

% ============================================================
\begin{frame}[fragile]{Stress test determinístico por clase de activo}
\begin{lstlisting}[style=Rstyle]
pesos <- c(bonos = 0.55, equity = 0.25, externo = 0.15, liquidez = 0.05)

shock_adverso <- c(
  bonos   = -0.12,
  equity  = -0.30,
  externo = -0.18,
  liquidez = 0.02
)

perdida_port <- sum(pesos * shock_adverso)
perdida_port
\end{lstlisting}

\begin{alertblock}{Interpretación}
La pérdida agregada se descompone por clase de activo, lo que facilita decisiones regulatorias sobre límites, concentración y cobertura.
\end{alertblock}
\end{frame}

% ============================================================
\begin{frame}[fragile]{Impacto sobre saldo acumulado}
\begin{lstlisting}[style=Rstyle]
saldo_inicial <- 100
r_base <- 0.05
r_adv  <- perdida_port

saldo_base <- saldo_inicial * (1 + r_base)
saldo_adv  <- saldo_inicial * (1 + r_adv)

impacto_saldo <- saldo_adv - saldo_base
impacto_pct   <- impacto_saldo / saldo_base

c(saldo_base = saldo_base,
  saldo_adverso = saldo_adv,
  impacto_pct = impacto_pct)
\end{lstlisting}

\begin{block}{Lectura previsional}
Un shock financiero transitorio puede tener efectos persistentes si ocurre cerca de la edad de retiro o si reduce la base sobre la cual se capitalizan retornos futuros.
\end{block}
\end{frame}

% ============================================================
\begin{frame}[fragile]{Correlación estresada: ejemplo simple}
\begin{lstlisting}[style=Rstyle]
vol <- c(bonos = 0.08, equity = 0.18, externo = 0.16)

R_normal <- matrix(c(
  1.00, 0.20, 0.30,
  0.20, 1.00, 0.55,
  0.30, 0.55, 1.00), 3, 3)

R_stress <- matrix(c(
  1.00, 0.65, 0.70,
  0.65, 1.00, 0.85,
  0.70, 0.85, 1.00), 3, 3)

Sigma_normal <- diag(vol) %*% R_normal %*% diag(vol)
Sigma_stress <- diag(vol) %*% R_stress %*% diag(vol)
\end{lstlisting}
\end{frame}

% ============================================================
\begin{frame}[fragile]{Comparación de volatilidad de portafolio}
\begin{lstlisting}[style=Rstyle]
w <- c(0.55, 0.25, 0.20)

vol_port_normal <- sqrt(t(w) %*% Sigma_normal %*% w)
vol_port_stress <- sqrt(t(w) %*% Sigma_stress %*% w)

c(normal = vol_port_normal,
  stress = vol_port_stress)
\end{lstlisting}

\begin{alertblock}{Mensaje}
El portafolio puede parecer diversificado con correlaciones normales, pero mucho menos resiliente bajo correlaciones de crisis.
\end{alertblock}
\end{frame}

% ============================================================
\section{Ejercicio}

\begin{frame}{Ejercicio aplicado: caso de fondo previsional}
\begin{block}{Portafolio inicial}
\begin{itemize}
  \item 55\% bonos soberanos y corporativos locales.
  \item 25\% acciones globales.
  \item 15\% renta fija externa.
  \item 5\% liquidez.
\end{itemize}
\end{block}

\begin{alertblock}{Escenario adverso proveniente del Módulo 3}
\begin{itemize}
  \item Aumento de tasas locales: +300 puntos básicos.
  \item Caída de equity global: -30\%.
  \item Aumento de spread soberano: +250 puntos básicos.
  \item Depreciación cambiaria: +18\%.
  \item Inflación anual: 9\%.
\end{itemize}
\end{alertblock}
\end{frame}

% ============================================================
\begin{frame}{Tareas del ejercicio}
\begin{enumerate}
  \item Traducir el escenario macro-financiero en retornos por clase de activo.
  \item Calcular pérdida total del portafolio.
  \item Descomponer contribución a pérdida por factor.
  \item Comparar baseline vs. adverso.
  \item Estimar VaR y Expected Shortfall.
  \item Evaluar impacto en saldo de un afiliado:
  \begin{itemize}
    \item joven: 30 años,
    \item medio: 45 años,
    \item próximo al retiro: 60 años.
  \end{itemize}
  \item Formular recomendaciones regulatorias.
\end{enumerate}
\end{frame}

% ============================================================
\begin{frame}{Plantilla de resultados}
\begin{table}
\centering
\scriptsize
\begin{tabular}{lrrrr}
\toprule
\textbf{Clase de activo} & \textbf{Peso} & \textbf{Ret. baseline} & \textbf{Ret. adverso} & \textbf{Contrib. pérdida} \\
\midrule
Bonos locales & 55\% & 6\% & -12\% &  \\
Acciones globales & 25\% & 8\% & -30\% &  \\
Renta fija externa & 15\% & 5\% & -18\% &  \\
Liquidez & 5\% & 3\% & 2\% &  \\
\midrule
\textbf{Total} & 100\% & & & \\
\bottomrule
\end{tabular}
\end{table}

\begin{block}{Pregunta de decisión}
¿Qué clase de activo explica la mayor pérdida y qué medida regulatoria sería proporcional al riesgo observado?
\end{block}
\end{frame}

% ============================================================
\begin{frame}{Discusión en clase}
\begin{columns}[T]
\column{0.48\textwidth}
\begin{block}{Preguntas técnicas}
\begin{itemize}
  \item ¿La pérdida viene de tasas, equity o FX?
  \item ¿Cuánto aumenta el riesgo con correlaciones de crisis?
  \item ¿Qué cambia si usamos t-Student en lugar de normal?
\end{itemize}
\end{block}

\column{0.48\textwidth}
\begin{block}{Preguntas regulatorias}
\begin{itemize}
  \item ¿Se requieren límites por duración?
  \item ¿Debe haber cobertura cambiaria mínima?
  \item ¿Conviene un glide path por edad?
  \item ¿Qué divulgar al afiliado?
\end{itemize}
\end{block}
\end{columns}
\end{frame}

% ============================================================
\section{Interpretación}

\begin{frame}{Lectura de resultados}
\begin{block}{No basta reportar pérdida total}
El reporte debe explicar:
\begin{itemize}
  \item qué factor explica la pérdida,
  \item si la pérdida es transitoria o persistente,
  \item qué grupo de afiliados es más vulnerable,
  \item qué margen de acción tiene el regulador,
  \item qué supuestos dominan los resultados.
\end{itemize}
\end{block}

\begin{alertblock}{Clave}
La interpretación convierte un ejercicio estadístico en una herramienta de supervisión.
\end{alertblock}
\end{frame}

% ============================================================
\begin{frame}{Implicaciones por cohorte de afiliados}
\begin{table}
\centering
\small
\begin{tabular}{p{0.25\textwidth} p{0.35\textwidth} p{0.30\textwidth}}
\toprule
\textbf{Cohorte} & \textbf{Vulnerabilidad} & \textbf{Respuesta de política} \\
\midrule
Jóvenes & Mayor tiempo para recuperar pérdidas & Portafolio más diversificado y crecimiento de largo plazo \\
Edad media & Riesgo de trayectoria y acumulación parcial & Rebalanceo y control de drawdowns \\
Próximos al retiro & Alta exposición al momento de conversión a pensión & Glide path, liquidez y menor duración \\
\bottomrule
\end{tabular}
\end{table}
\end{frame}

% ============================================================
\begin{frame}{Decisiones regulatorias vinculadas al Módulo 4}
\begin{itemize}
  \item Límites de concentración por emisor, sector y moneda.
  \item Límites de duración o sensibilidad a tasas.
  \item Reglas de cobertura cambiaria.
  \item Requerimientos de liquidez para retiros, traspasos o beneficios.
  \item Escenarios regulatorios comunes para comparar AFP.
  \item Reporte de pérdidas por factor y por cohorte.
  \item Revisión de multifondos y reglas de ciclo de vida.
\end{itemize}
\end{frame}

% ============================================================
\begin{frame}{Formato mínimo de reporte regulatorio}
\begin{enumerate}
  \item \textbf{Escenario}: baseline y adverso, con narrativa y supuestos.
  \item \textbf{Metodología}: modelo multifactorial, GARCH, correlaciones y simulación.
  \item \textbf{Resultados}: pérdida total, VaR, ES, drawdown y saldo.
  \item \textbf{Descomposición}: contribución por activo, factor y moneda.
  \item \textbf{Sensibilidades}: tasas, FX, equity, spreads e inflación.
  \item \textbf{Recomendaciones}: límites, coberturas, liquidez y divulgación.
\end{enumerate}
\end{frame}

% ============================================================
\begin{frame}{Mensajes clave}
\begin{enumerate}
  \item El Módulo 3 produce trayectorias macro-financieras; el Módulo 4 las transforma en pérdidas de portafolio.
  \item El riesgo financiero es riesgo previsional cuando afecta el saldo acumulado.
  \item Las betas muestran los canales de transmisión.
  \item GARCH permite capturar aumento y persistencia de volatilidad.
  \item Las correlaciones de crisis reducen la diversificación efectiva.
  \item VaR y ES deben complementarse con escenarios narrativos plausibles.
  \item El resultado final debe informar decisiones regulatorias.
\end{enumerate}
\end{frame}



% ============================================================
\begin{frame}
\centering
\Huge Gracias
\vspace{0.6cm}


\end{frame}

\end{document}
